\begin{table*}[htbp]
\centering
\scriptsize
\caption{Классификатор: качество по вариантам модели и ablation структурных признаков}
\label{tab:4}
\begin{tabular}{>{\raggedright\arraybackslash}p{0.115\linewidth}>{\raggedright\arraybackslash}p{0.115\linewidth}>{\raggedright\arraybackslash}p{0.115\linewidth}>{\raggedright\arraybackslash}p{0.115\linewidth}>{\raggedright\arraybackslash}p{0.115\linewidth}>{\raggedright\arraybackslash}p{0.115\linewidth}>{\raggedright\arraybackslash}p{0.115\linewidth}>{\raggedright\arraybackslash}p{0.115\linewidth}}
\toprule
Вариант модели & Что входит в признаки & Holdout с определённым AUROC & Медиана AUROC & Размах AUROC & Медиана FPR & Контраст P3 − P1, полная матрица & Контраст P3 − P1, без структурных признаков \\
\midrule
source-only & только идентификатор площадки & 13 из 18 & 1.0000 & 0.5000–1.0000 & 1.0000 & — & — \\
main+M02 & плюс признак сходства абзацев & 13 из 18 & 0.9979 & 0.9606–1.0000 & 0.0332 & — & — \\
main & 22 признака стиля & 13 из 18 & 0.9973 & 0.9433–1.0000 & 0.0380 & −0.0941 [−0.1473; −0.0463], p = 0.0028 & +0.0169 [−0.0134; 0.0496], p = 0.3168 \\
format-only & четыре структурных признака & 13 из 18 & 0.9775 & 0.9154–1.0000 & 0.0118 & — & — \\
genre-only & только жанр & 13 из 18 & 0.6667 & 0.5000–0.7259 & 1.0000 & — & — \\
length-only & только длина текста & 13 из 18 & 0.4364 & 0.2957–0.6138 & 0.6276 & — & — \\
negative-control & перемешанные метки классов & 13 из 13 & 0.3614 & 0.1279–0.4440 & — & — & — \\
\bottomrule
\end{tabular}
\begin{minipage}{\linewidth}\footnotesize\vspace{2pt}
\textit{Примечание.} Единица анализа — метрики качества — групповой holdout; контрасты — пара текстов, порождённых одной моделью по одному заданию. Знаменатель: 18 групповых holdout; контраст считается на 120 парах при 15 кластерах, только на жанре seo. Данные: серия v2, статус current. Медиана и размах по holdout; контраст — абсолютная разница вероятностей с 95\% интервалом и p дикого кластерного бутстрапа. Интервалы 95\%, кластерный бутстрап по автору и документу; кластер — задание, заданий 45. Шкала: AUROC: больше значит лучше разделение; FPR: меньше значит меньше ложных обвинений человека; контраст в вероятности класса A: больше значит более AI-подобно. Сокращения: AUROC — площадь под ROC-кривой; FPR — доля ложных срабатываний на человеческих текстах; ablation — сравнение модели с признаками и без них, не разложение балла на слагаемые. Пропуски: AUROC не определён на пяти одноклассовых holdout, поэтому в колонке указано, на скольких он посчитан; контраст считался только для основной модели, у прочих вариантов стоит прочерк; у модели negative-control FPR не считался, а holdout тринадцать: она прогонялась серией кластерных перестановок и получила статус post hoc. У модели source-only AUROC 1.0000 сочетается с FPR 1.0000: ранжирование по площадке безошибочно, но при пороге 0.5 модель объявляет машинным весь человеческий тест невиданной площадки. Семейство первичных исходов закрыто поправкой Бонферрони на два зарегистрированных теста, α = 0.025 на тест; интервалы нескорректированные.
\end{minipage}
\end{table*}
